% Background on Software Defined Networking.

\section{Software Defined Networking}
\label{sec:sdn}

Software defined networking (SDN)~\cite{10.1145/1355734.1355746} decouples the
\emph{control plane} (routing decisions, topology policy, admission rules) from
the \emph{data plane} (packet forwarding on switches and routers).
In a traditional network each device runs distributed control software and
forwards according to locally derived state; in SDN, a logically centralized
\emph{controller} maintains a global view of the network and programs
forwarding tables on the data plane through a southbound interface such as
OpenFlow~\cite{10.1145/1355734.1355746}.
The data plane remains responsible only for fast, deterministic packet
processing (match, forward, enqueue), while the control plane can evolve
independently, be updated without changing hardware, and apply consistent
policy across many nodes in one step.
That separation lowers operational cost, enables programmatic automation, and
allows researchers and operators to experiment with new routing and management
logic without replacing every network element.

A natural way to give the controller that global view is to represent the
network as a \emph{graph}: vertices are hosts, switches, and links; edges
encode connectivity and weights such as delay or capacity.
Graphs provide a unified structure on which to run management algorithms
(shortest path, spanning tree, coloring) and to notify applications when
topology changes~\cite{7014202}.
OrbitGraph (Section~\ref{sec:orbitgraph}) follows this graph-centric SDN
principle for LEO constellations: the controller builds a time-varying graph
$G_t$ from orbital snapshots, computes forwarding centrally, and pushes only
the deltas required to keep the data plane aligned with each graph edit.
